\section{Matériaux pour les batteries}

\setcounter{equation}{0}
\setcounter{Q}{0}

Les batteries lithium-ion sont les batteries que nous utilisons dans tous nos objets du quotidien. Dans ce TD, nous allons nous intéresser au fonctionnement de ces batteries et au rôle du chimiste des matériaux pour leur conception. Dans la Figure 15 ci-dessous, vous trouverez les différents composants d’une batterie lithium ion. Dans l’ensemble du TD, nous garderons le même matériau à l’anode : le graphite. Il s’agit d’un allotrope du carbone couramment utilisé dans lequel le lithium vient s’intercaler entre les feuillets 2D de graphène. Cette anode sera notée \ce{LiC6}. 
La cathode est le matériau qui a fait l’objet du plus de développement en recherche pendant ces 20 dernières années. Il existe trois grandes familles pour ces matériaux de batteries et nous en aborderons majoritairement deux : les oxydes lamellaires et les structures de type spinelle. 

\begin{center}
\includegraphics[width=10cm]{Ex06/Fig01}
\captionof{figure}{Composants d’une batterie Lithium-ion}
\end{center}

\subsection*{Considérations générales} 

Tout d’abord, pour évaluer les performances d’une batterie, il est nécessaire de prendre en compte plusieurs paramètres :
\begin{itemize}
\item La tension de la batterie. Plus la tension est élevée, plus la puissance de la batterie est importante.
\item La cyclabilité : il s’agit grossièrement du temps de vie de la batterie.
\item La tenue en température de la batterie. 
\item Sa capacité spécifique (mesurée en Ah/kg) qui correspond à la quantité de courant présent dans la batterie par unité de masse. Plus cette capacité est élevée, plus la batterie pourra délivrer du courant (elle aura donc une meilleure autonomie). La formule pour calculer la capacité spécifique est donnée ci-dessous. 
\end{itemize}

\begin{equation}
C = 26.8 \frac{\text{nombre de lacune}}{m}
\end{equation}

Où $m$ est la masse de la batterie considérée (\si{\kilo\gram}) ; le nombre de lacunes correspond à la quantité de charges (électrons ou lithium) qui peuvent être stockés dans la batterie et le facteur de $26.8$ est un facteur lié à la conversion des Coulomb vers les Ampère heure. 

\Q Pourquoi les matériaux denses ne conviennent pas pour élaborer des batteries ?
\Q Sur le schéma de la Figure 15, indiquer dans quels sens circulent les électrons et les ions lithium durant la charge et la décharge. 
\Q Ecrire les demi-équations électroniques associées. Vous noterez le matériau cathodique comme LiMO2.  

Comment choisir le métal M ? 
Un des premiers points à considérer est l’ionicité de la liaison métal oxygène. En effet, plus la liaison est ionique, plus le potentiel d’oxydoréduction associé augmente. 
\Q Rappeler l’évolution de l’électronégativité au sein du tableau périodique. En déduire pourquoi il est préférable de travailler avec des oxydes plutôt que des sulfures ou des sélénures. 
\Q Par quoi peut-on remplacer l’oxygène pour avoir des matériaux avec une liaison M-X encore plus ionique. Vous pouvez proposer plusieurs hypothèses.

Un autre paramètre clé est l’arrangement des polyèdres en 3D. Ils peuvent soit se partager des faces soit des arrêtes, soit des sommets (voir Figure 16 ci-dessous). Cet arrangement joue sur la covalence des liaisons metal/oxygène et la covalence augmente en passant de polyèdres partageants des faces et des arrêtes (\ce{Fe2O3}) à des polyèdres partageant des sommets (ex : \ce{Na4Fe2O5}).

\begin{center}
\includegraphics[width=10cm]{Ex06/Fig02}
\captionof{figure}{Exemples d’arrangements de polyèdres différents dans des différentes phases d’oxyde de fer.} 
\end{center}

\Q Expliquer pourquoi l’augmentation du caractère covalent des liaisons metal/oxygène ne nous arrange pas pour obtenir des matériaux performants. 

\subsection*{L’oxyde lamellaire \ce{LiCoO2}}

La première batterie commercialisée par Sony en 1990 comportait l’oxyde lamellaire LiCoO2 dont la structure est donnée sur la Figure 17.
\Q Quelle est la charge électrique des feuillets composés de polyèdres cobalt-oxygène ? 
\Q Cette première batterie possédait une capacité spécifique de 150 mAh/g ce qui est relativement faible (on peut maintenant obtenir le double). Proposer une hypothèse expliquant cette relativement faible capacité. 



\begin{center}
\includegraphics[width=10cm]{Ex06/Fig03}
\captionof{figure}{Représentation de la structure lamellaire de l’oxyde \ce{LiCoO2}}
\end{center}

\Q Comment peut-on améliorer la situation ? Pourquoi ? 

\subsection*{Les matériaux de type spinelle (\ce{AB2X4})}


Certains spinelles peuvent être utilisés dans le domaine des batteries. En effet, ce type de matériaux présente une structure en 3D qui permet de capter les ions lithium comme montré dans la figure ci-dessous pour la strucuture spinelle LixMn2O4. 

\Q Rappeler la différence entre spinelle normale et spinelle inverse.

Données :
$r(O2-) = 140 pm ; r(Li+) = 76 pm ; r(Mn3+) = 67 pm$

\begin{center}
\includegraphics[width=10cm]{Ex06/Fig04}
\captionof{figure}{Projection 2D de la structure \ce{LiMn2O4} (a) et structure 3D (b)}
\end{center}

\Q Donner le degré d’oxydation des atomes de manganèse dans la composition limite LiMn2O4. 
A l’état solide, l’oxyde de lithium manganèse cristallise dans un réseau de type spinelle, où les ions \ce{O2-} occupent les n{\oe}uds d’un réseau cubique à faces centrées, les ions métalliques se répartissant sur certains sites octaédriques et tétraédriques du réseau.
\Q Indiquer la position des sites octaédriques et tétraédriques dans le réseau cubique à faces centrées. Dénombrer ces sites. 
\Q Evaluer le rayon maximal $r$ d’un atome sphérique pouvant se loger sans déformation du réseau CFC, au sein d’un site tétraédrique. 
\Q Compte tenu des rayons ioniques, la structure est-elle déformée ? 
\Q Sur des critères purement géométriques, la structure doit-elle être directe ou inverse ?
\Q Expliquer pourquoi elle cristallise selon une structure de type spinelle directe. 

Le composé \ce{LixMn2O4} a une durée vie relativement faible car le manganèse diffuse dans la batterie et a des performances relativement limitées. En effet, lorsque x = 1, il y a autant de \ce{Mn3+} dans la structure que de \ce{Mn4+} et cela provoque une chute énorme de potentiel (donc la batterie ne délivre plus de puissance). Nous allons essayer de comprendre pourquoi. 

\Q Donner la configuration électronique d’un ion \ce{Mn3+}. Comment sont répartis les électrons dans les orbitales $3d$ ?
\Q Du fait de cette répartition, les complexes octaédriques à base de Mn3+ sont distordus. Cet effet, appelé Jahn Teller, provoque une augmentation du volume de l’octaèdre. Expliquer pourquoi cette expansion est problématique dans le cas d’une batterie. 
1
\Q Comment pouvons-nous remédier à ce problème ? 

\subsection*{Conclusion :}
Aujourd’hui, des structures modifiées et raisonnées sont utilisées notamment en tant que batterie de voiture. Par exemple le composé \ce{Li Mn_{2-y} Al_y O_{4-z} F_z}. 
La recherche sur les batteries est toujours riche et nécessite encore des années de travaux. En effet, certaines difficultés doivent encore être levées :
\begin{itemize}
\item Réduire l’empreinte carbone des batteries en choisissant des métaux abondants (ou recyclés), en développant des méthodes de synthèse de chimie douce.
\item Développer des batteries tout solides pour éviter les électrolytes liquides ou travailler en milieu aqueux pour améliorer la sécurité.
\item Substituer le lithium par du sodium, du calcium ou encore du magnésium. 
\item Développer d’autres technologies de batterie :
\end{itemize}


\begin{center}
\includegraphics[width=10cm]{Ex06/Fig05}
\captionof{figure}{Différentes technologies de batteries en cours de développement}
\end{center}
